
\documentclass[aspectratio=169]{beamer}

\mode<presentation>
\usetheme{Boadilla}
\definecolor{NTHU1}{RGB}{127,16,132}
\definecolor{NTHU2}{RGB}{228,210,231}
\definecolor{blue}{RGB}{30,90,205}
\definecolor{red}{RGB}{213,94,0}
\definecolor{green}{RGB}{0,128,0}
\definecolor{charcoalgray}{RGB}{108,104,104}
\setbeamercolor{title}{fg=NTHU1}
\setbeamercolor{frametitle}{fg=NTHU1}
\setbeamercolor{block title}{bg=NTHU1, fg=white}
\setbeamercolor{block body}{bg=NTHU2}
\setbeamercolor{structure}{fg=NTHU1}
\setbeamercolor{item projected}{fg=white}
\setbeamercolor{item}{fg=NTHU1}
\setbeamercolor{subitem}{fg=NTHU1}
\setbeamercolor{section in toc}{fg=NTHU1}
\setbeamercolor{description item}{fg=NTHU1}
\setbeamercolor{caption name}{fg=NTHU1}
\setbeamercolor{button}{bg=NTHU1, fg=white}
\setbeamercolor{caption name}{fg=NTHU1}
\usepackage{graphics}
\usepackage{tikz}
\usepackage{amsmath}
\usepackage{bbm}
\usetikzlibrary{decorations.pathreplacing}
\usepackage{geometry}
\usepackage{booktabs}
\usepackage{bookmark}
\usepackage{multirow, makecell}
\usepackage{float}
\usepackage{fancyvrb}
\usepackage{caption}
\captionsetup[figure]{singlelinecheck = false, format= hang, justification=centering}
\captionsetup[subfigure]{singlelinecheck = false, format= hang, justification=raggedright}
\usepackage{subcaption}
\usepackage{adjustbox}
\usepackage{hyperref}
\usepackage{threeparttable}
\usepackage{appendixnumberbeamer} 
\usepackage[T1]{fontenc}
\usepackage[default]{lato} 
\usepackage{smartdiagram}
\smartdiagramset{
  back arrow disabled = true,
  module minimum width=1cm,
  module x sep = 3,
  uniform arrow color=true,
}
\usepackage{xcolor}
\usepackage{colortbl}
  \definecolor{light-gray}{gray}{0.99}

\newenvironment{wideitemize}{\itemize\addtolength{\itemsep}{10pt}}{\enditemize}
\newenvironment{wideenumerate}{\enumerate\addtolength{\itemsep}{10pt}}{\endenumerate}
\newenvironment{widedescription}{\description\addtolength{\itemsep}{10pt}}{\enddescription}
\hypersetup{
colorlinks=true,
linkcolor=NTHU1,
filecolor=green, 
urlcolor=blue,
}
\beamertemplatenavigationsymbolsempty
\setbeamercolor{author in head/foot}{bg=white, fg=NTHU1}
\setbeamercolor{title in head/foot}{bg=white, fg=NTHU1}
\setbeamercolor{date in head/foot}{bg=white, fg=NTHU1}
\setbeamercolor{section in head/foot}{bg=white, fg=NTHU1}
\setbeamercolor{headline}{bg=white}
\setbeamertemplate{footline}{
    \leavevmode%
    \hbox{%
        \begin{beamercolorbox}[wd=.333333\paperwidth,ht=2.25ex,dp=1ex,center]{date in head/foot}%
            \usebeamerfont{date in head/foot}%\insertdate
        \end{beamercolorbox}%
        \begin{beamercolorbox}[wd=.333333\paperwidth,ht=2.25ex,dp=1ex,center]{title in head/foot}%
            \usebeamerfont{title in head/foot}\insertshorttitle
        \end{beamercolorbox}%
        \begin{beamercolorbox}[wd=.333333\paperwidth,ht=2.25ex,dp=1ex,right]{date in head/foot}%
            \usebeamerfont{date in head/foot}\insertframenumber{}\hspace*{2em}
        \end{beamercolorbox}}%
        \vskip0pt%
    }
    %% May need to script the introduction!!!!!!
%\setbeamercolor{page number in head/foot}{fg=NTHU1}
\setbeamertemplate{section in toc}[sections numbered]
\setbeamertemplate{subsection in toc}{\leavevmode\leftskip=3em\rlap{\hskip-1.75em\inserttocsectionnumber.\inserttocsubsectionnumber}
\inserttocsubsection\par}
\setbeamerfont{subsection in toc}{size=\footnotesize}
%\setbeamertemplate{headline}{%
  %\begin{beamercolorbox}[ht=5.5ex]{section in head/foot}
    %\vskip2pt\insertnavigation{0.33\paperwidth}\vskip2pt
  %\end{beamercolorbox}%
%}
\newenvironment{transitionframe}{\setbeamercolor{background canvas}{bg=white}\setbeamertemplate{footline}{} \begin{frame}[noframenumbering]}{\end{frame}}

\makeatletter
\let\@@magyar@captionfix\relax
\makeatother

\title[]{Conflict and Peace} % Change this regularly
\subtitle[]{Week 3: Ethnic and Political Origins of Conflicts}
\author[Lee]{Seung-hun Lee }
\institute[]{Taipei School of Economics\\ National Tsing Hua University}

\date[]{September 19th, 2025}

\begin{document}
\begin{transitionframe}
\titlepage
\end{transitionframe}

%%% Color slides for section headers: Use for colloquium version (The ones bewteen \iffals and \fi)

\section{Overview of the topic}
\begin{frame}
\frametitle{How do institutional backgrounds shape conflicts?} 
Institutional backgrounds can interact with economic factors that affect conflicts
\begin{itemize}\setlength\itemsep{3pt}
\item Democracies can hold decisionmakers accountable
\item Transparency reduces uncertainties and misperceptions
\item Cohesive institutions can defuse conflict by bridging different groups/stakeholders
\end{itemize}\medskip
Yet, there are many cases where institutions can foster conflict
\begin{itemize}\setlength\itemsep{3pt}
\item Many states are still autocratic
\item Many democracies are also experiencing surge of extremist movements
\item Media are often misused as tools of propaganda and misinformation, not transparency
\item Leads to difficulty in fostering cohesion: Cultural and ethnic wars on the rise
\end{itemize}
\end{frame}

\begin{frame}
\frametitle{Conflict across countries with different political spectrum}
\begin{figure}
  \centering
  \includegraphics[keepaspectratio,width=0.8\textwidth]{fig1.png}
\end{figure}
\end{frame}

\begin{frame}
\frametitle{Political motives and violence}
\begin{figure}
  \centering
  \includegraphics[keepaspectratio,width=0.9\textwidth]{fig2.png}
\end{figure}
\end{frame}

\begin{frame}
\frametitle{Where we are at} 
Need to understand institutions and economics behind conflicts
\begin{itemize}\setlength\itemsep{3pt}
\item Context-by-context: Fractionalized and polarized ethnicities, religion
\item Attempt to understand what institutional backgrounds shapes economic motives
\item Possible with innovative research design based on new numerical, satellite, text data 
\end{itemize}\medskip

We explore the cultural origins of war: Ethnicity, colonialism, religion etc \\
\medskip
We also look at how political motives shape wars: Propaganda, misinformation, strategic moves for political gain in elections etc.
\end{frame}

\begin{transitionframe}
\frametitle{Roadmap for today}
\tableofcontents
\end{transitionframe}

\section{Institutional and cultural origins of conflict}

\subsection{Esteban, Mayoral, and Ray (2012): Ethnicity and Conflict: an Empirical Study}
\begin{transitionframe}
\begin{center}
  \begin{huge}
    \textcolor{NTHU1}{%
Esteban, Mayoral, and Ray (2012): \\ Ethnicity and Conflict: an Empirical Study
    }
  \end{huge}
\end{center}
\end{transitionframe}

\begin{frame}
\frametitle{Q: How do ethnic divisions affect conflict?} 
Are there meaningful relationship between ethnic division and social conflict?
\begin{itemize}\setlength\itemsep{3pt}
\item Early on: Income and wealth distribution $\to$ conflict?
\begin{itemize}
\item Poor antagonistic to the rich
\end{itemize}
\item Lacked explanation on how this sentiment channels into organized action
\item Start looking into different markers: Identities, ethnicities, religions etc
\end{itemize} \medskip
This paper extends the explanation of sources of conflicts beyond economic factors
\begin{itemize}\setlength\itemsep{3pt}
\item In particular, focuses on how fractionalized the society is along ethnic lines
\item Provides definition of ethnic division and how each relate to conflicts
\item Heterogeneities depending on conflict over public goods and private goods 
\begin{itemize}\setlength\itemsep{3pt}
  \item Public: Benefits independent of group size, but depends on identity (ethnic dominance)
  \item Private goods: Benefits decrease with group size, but identity does not matter (resources)
  \end{itemize}
\item Empirically demonstrates their finding at a country level data
\end{itemize} \medskip
\end{frame}

\begin{frame}
\frametitle{Setting up the ethnic conflict model} 
Setup: Total of $m$ groups within a society in conflict over public and private goods
\begin{itemize}\setlength\itemsep{3pt}
\item Group $i$ has $N_i$ people (Total population: $N$), with $n_i$ share of population
\item Each individual exerts effort $r$, with cost of that effort being $c(r)$
\item Each group exerts $R_i$ effort, with probability of winning being $p_i=\frac{R_i}{R_1+...+R_N}$
\item Private goods of $\mu$  and public good of $\pi$ at stake (per capita)
\end{itemize} \medskip
What are the payoffs of conflict for group $i$?
\begin{itemize}\setlength\itemsep{3pt}
  \item Private goods are divided among same groups
  \item If group $i$'s desired public good chosen, $\pi u_{ii}$, otherwise incurs cost $\pi(u_{ii}-u_{ij})=\pi d_{ij}$
  \[
\underbrace{\pi u_{ii} -\sum_{j=1}^m p_j \pi d_{ij}}_{\text{Net gains of public good $j$ chosen for $i$}} + \underbrace{p_i \frac{\mu}{n_i}}_{\text{Gains from private good}} - \underbrace{c(r)}_{\text{cost of effort}}
  \]
\end{itemize}
\end{frame}

\begin{frame}
\frametitle{Definition of indices of ethnic division} 
Different types of measures
\begin{itemize}\setlength\itemsep{3pt}
\item Polarization: Emphasis of number \textit{and symmetry in distribution} of groups in conflict
\[
P = \sum_{i=1}^m\sum_{j=1}^m n_i^2 n_j d_{ij}
\]
\item Greenberg-Gini Index: Emphasis on the number of groups in conflict
\[
G = \sum_{i=1}^m\sum_{j=1}^m n_i n_j d_{ij}
\] 
\item Fractionalization: Generalized G to case where $d_{ii}=0$ and $d_{ij}=1$ for $j\neq i$
 \[
F = \sum_{i=1}^m n_i(1- n_i)
\]

\end{itemize} \medskip
\end{frame}


\begin{frame}
\frametitle{Hypotheses and data} 
Hypotheses
\begin{itemize}\setlength\itemsep{3pt}
\item Public goods: Polarization increases conflicts (symmetric groups $\to$ contention $\uparrow$)
\item Private goods: Fractionalization increases conflict  (Division of limited goods)
\item Group cohesion (low $G/N$, groups per capita) increases conflict (easier mobilization)
\end{itemize} \medskip
Dataset
\begin{itemize}\setlength\itemsep{3pt}
\item Conflicts: UCDP/PRIO dataset 
\item Definition of groups: Ethnolinguistic dataset from Fearon (2003), among others 
\item Also control for GDP, population, resources (oil), terrains, democracy, past conflict
\item Runs OLS to map correlation between conflict and indices
 \[
Y_{it} = \alpha + \beta_1 P + \beta_2 F + \beta_3 (G/N) + \gamma X_{it} + e_{it}
 \]
\end{itemize} \medskip
\end{frame}

\begin{frame}
\frametitle{Baseline findings} 
\begin{figure}
  \includegraphics[keepaspectratio,width=0.9\textwidth]{emr2012_t1.jpeg}
\end{figure}
\begin{itemize}\setlength\itemsep{3pt}
\item p-values in parenthesis
\item Polarization explains conflicts far more than Fractionalization
\item Group cohesion (smaller G/N) increases conflicts
\end{itemize}
\end{frame}

\begin{frame}
\frametitle{Graphical findings} 
\begin{figure}
  \includegraphics[keepaspectratio,width=0.9\textwidth]{emr2012_f1.jpeg}
\end{figure}
\begin{itemize}\setlength\itemsep{3pt}
\item Polarization has steeper slope
\end{itemize}
\end{frame}

\begin{frame}
\frametitle{Discussion and takeaways} 
Takeaways
\begin{itemize}\setlength\itemsep{3pt}
  \item Each conflict contains fight for private and public goods
  \item Polarization explains the incidence of conflicts, especially those with public nature
  \item Fractionalization explains conflicts too, albeit with limited scope
  \item Group cohesion increases conflicts
\end{itemize}\medskip
Limits and going forward
\begin{itemize}\setlength\itemsep{3pt}
  \item Empirically separating private and public component: Not mutually exclusive
  \item Applicable to other divides: Religion, caste, language etc
  
\end{itemize}\medskip
\end{frame}

\subsection{Bonomi, Gennaioli, Tabellini (2021): Identity, Beliefs, and Political Conflict}
\begin{transitionframe}
\begin{center}
  \begin{huge}
    \textcolor{NTHU1}{%
Bonomi, Gennaioli, Tabellini (2021): \\ Identity, Beliefs, and Political Conflict
    }
  \end{huge}
\end{center}
\end{transitionframe}

\begin{frame}
\frametitle{Q: How does identity affect behavior of groups in conflict} 
Increasing salience of conflicts over identity and culture
\begin{itemize}\setlength\itemsep{3pt}
\item Conflict (not necessarily violent) over immigration, race, and ethnicity on the rise
\begin{itemize}
\item Preference for redistribution increasingly correlated with progressive cultural policy
\end{itemize}
\item Growing partisanship and hatred towards outsiders
\item While mostly confined to politics, it can extend to violent conflicts
\end{itemize}\medskip
This paper models how identity becomes important in forming beliefs
\begin{itemize}\setlength\itemsep{3pt}
\item Cultural conflict $\Uparrow$ $\to$ salience of identity $\Uparrow$ $\to$ Belief on group-linked stereotypes $\Uparrow$
\item If culture becomes more relevant for welfare, it becomes key determinant of identity
\item Similar if economic shocks affects cultural groups differently (globalization)
\item Increased belief distortions leading to rise of ingroup vs outgroup conflict
\end{itemize}\medskip
\end{frame}

\begin{frame}
\frametitle{Model of identity and political conflict} 
Setup of policies and voters
\begin{itemize}\setlength\itemsep{3pt}
\item Policy 1: Proportional income tax $\tau$ that finances public good, but has distortion $-\frac{\varphi}{2}\tau^2$
\item Policy 2: Cultural policy $q$ (race, immigration) where $q\Uparrow\to$ more liberal stance
\item Voters are defined by relative income $\epsilon$ and cultural trait $\psi$ (mean 0, st.dev 1)
\begin{itemize}\setlength\itemsep{3pt}
\item Higher $\epsilon$: Relatively wealthy; Higher $\psi$: Prefer progressive values
\item $\rho$ captures correlation between $\epsilon$ and $\psi$: We assume non-negative correlation
\end{itemize}
\item Public goods have marginal value $\nu$, and $\psi$ types value it additionally by $\beta$
\end{itemize}\medskip
Policy preference is the sum of expected income, public goods, and cultural policy
\[
W^{\epsilon\psi}(\tau,q)=\underbrace{(1+\epsilon)(1-\tau)-\frac{\varphi}{2}\tau^2}_{\text{Income net of taxes and distortions}}+\underbrace{(\nu+\beta\psi)\tau}_{\text{Value of public goods to $\psi$}}-\underbrace{\frac{\kappa}{2}(q-\psi)^2}_{\text{preference on $q$, weighed by $\kappa$}}
\]
\end{frame}

\begin{frame}
\frametitle{Optimal policy preference} 

How do voters evaluate policies?
\[
W^{\epsilon\psi}(\tau,q)=\underbrace{(1+\epsilon)(1-\tau)-\frac{\varphi}{2}\tau^2}_{\text{Income net of taxes and distortions}}+\underbrace{(\nu+\beta\psi)\tau}_{\text{Value of public goods to $\psi$}}-\underbrace{\frac{\kappa}{2}(q-\psi)^2}_{\text{preference on $q$, weighed by $\kappa$}}
\]
\begin{itemize}\setlength\itemsep{3pt}
\item The optimal individual (involves calculs) and social (plug in mean values) policies are 
\begin{itemize}\setlength\itemsep{3pt}
\item Individual: $\tau=\frac{(\nu+\beta\psi)-(1+\epsilon)}{\varphi}, q=\psi$; \ \ \ \  \ \ \    Social: $\tau=\frac{\nu-1}{\varphi}, q=0$
\item More progressive ($\psi\uparrow$): More redistribution and more liberal; Richer ($\epsilon\uparrow$): Less tax
\item Social optimum reflects average income and identities
\end{itemize}
\end{itemize}\medskip
Identities change, especially if only one categorization defines an identity
\begin{itemize}\setlength\itemsep{3pt}
\item Categories: Wealth (poor or rich) and culture (progressive or conservative)
\item Identity is set to maximize similarities within groups and conflicts between groups
\item Culture matters with high $\kappa$, inter-cultural differences, and $\rho$
\end{itemize}
\end{frame}

\begin{frame}
\frametitle{How changes in identity affects policy preferences?} 
Changes in salience of identities distort beliefs that decide preferences for $\tau$ and $q$
\begin{itemize}\setlength\itemsep{3pt}
\item Belief: One's perceived future income and cultural trait relative to real mean values
\[
\epsilon^\theta_G=\epsilon+\theta(\bar{\epsilon}_G - \bar{\epsilon}_{\text{not}\ G}), \ \ \ \  \psi^\theta_G=\psi+\theta(\bar{\psi}_G - \bar{\psi}_{\text{not}\ G})
\]
\item Distortion comes from perceived differences between in-group and out-groups
\end{itemize}\medskip
Distortion of beliefs $\to$ Distortion of preferred policies
\begin{itemize}\setlength\itemsep{3pt}
\item Poor and conservative voter: Demand for redistribution $\Uparrow$ ($\tau\Uparrow$) \& Liberal policy $\Downarrow$($q\Downarrow$)
  \item Salience of conservatism $\Uparrow$ ($\psi\Downarrow$) $\to$  preference for $q$ $\Downarrow$ \& $\tau$ $\Downarrow$
\begin{itemize}
\item Salience of culture over income decreases concerns about being poor (more optimistic $\epsilon$)
\end{itemize}
\item This could be triggered by shocks irrelevant to redistribution (e.g. migrant inflow)
\item Explain why some group of voters change parties entirely
\begin{itemize}
\item Conservative and poor voters: Left-wing (class salient) $\to$ Right-wing (culture salient)
\end{itemize}
\end{itemize}\medskip
\end{frame}

\begin{frame}
\frametitle{Salience of immigration shifted policy preferences} 
Salience of immigration $\Uparrow$ (early 2010s) $\to$ Shift policy preferences across different groups
\begin{center}
\begin{figure}
  
\begin{subfigure}{0.475\textwidth}
  
\includegraphics[keepaspectratio, width=0.9\textwidth]{bgt2021_f1.jpeg}
\caption{Salience of issues over time}
\end{subfigure}
\hfill
\begin{subfigure}{0.475\textwidth}
  
\includegraphics[keepaspectratio, width=0.725\textwidth]{bgt2021_f2.jpeg}
\caption{Group conflict over time}
\end{subfigure}
\end{figure}
\end{center}
\end{frame}

\begin{frame}
\frametitle{Discussion and takeaways} 
Takeaways from the paper
\begin{itemize}\setlength\itemsep{3pt}
\item Voters decide based on the perspective of the most salient group (income vs culture)
\item Changes in salience of certain groups $\to$ Changes in preferred policies 
\item These could change what is being contested in politics (and beyond)
\end{itemize}\medskip
Potential extensions
\begin{itemize}\setlength\itemsep{3pt}
\item Changes in salience of identity explains economic outcomes
\begin{itemize}
\item How does migration affect tax preferences? \\ (Alesina, Miano, Stantcheva 2023, to be covered later in course)
\end{itemize}
\item Violent conflicts could explained by this framework
\begin{itemize}
\item If what is at stake is resources and vengeance, conflicts can happen with wars
\end{itemize}
\item Potential implications for understanding modern party politics
\end{itemize}\medskip
\end{frame}

\subsection{Michalopoulos, Papaiaonnou (2016): The Long-Run Effects of the Scramble for Africa}
\begin{transitionframe}
\begin{center}
  \begin{huge}
    \textcolor{NTHU1}{%
Michalopoulos, Papaiaonnou (2016): \\ The Long-Run Effects of the Scramble for Africa
    }
  \end{huge}
\end{center}
\end{transitionframe}

\begin{frame}
\frametitle{Q: How does colonial border design affect conflict and development?} 
Long-lasting consequences of colonial institutions
\begin{itemize}\setlength\itemsep{3pt}
\item Extractive institutions $\to$ long-run impacts on economic \& human capital development
\item Less highlighted: Design of arbitrary borders without local knowledge
\begin{itemize}\setlength\itemsep{3pt}
\item Territorial borders became strict (not the case before)  
\item Many ethnic groups were spilt to multiple countries
\item Many countries ended up with multiple ethnic groups
\end{itemize}
\end{itemize}\medskip
This paper analyzes the long-run impact of such colonial institution
\begin{itemize}\setlength\itemsep{3pt}
\item Compare partitioned ethnic groups to non-split groups
\item Do they have different characteristics (geography, pre-colonial economics etc)?
\item Are partitioned groups susceptible to conflicts? What about economic development?
\end{itemize}
\end{frame}

\begin{frame}
\frametitle{Scramble for Africa} 
Timeline
\begin{itemize}\setlength\itemsep{3pt}
\item 1860s: French and British enter West Africa under agreement on areas of influence
\item More Europeans enter $\to$ Divide unexplored areas (e.g.1885 Berlin Conference)
\item Inadequate local knowledge (no African representatives)
\end{itemize} \medskip
Who are the partitioned ethnic groups?
\begin{itemize}\setlength\itemsep{3pt}
\item Ethnicities with 10\%+ of their surface belonging to multiple countries
\item 229 ethnicities (27.7\% of sample)
\end{itemize} \medskip
How partitioning can affect conflicts and development
\begin{itemize}\setlength\itemsep{3pt}
\item Secessionism: 20\% of African civil wars involve secessionist component
\item Repression: Ethnicity-based discrimination on economic and social activities
\item Spillover: Population displacement $\to$ Ethnic component change in neighbors
\end{itemize} \medskip
\end{frame}

\begin{frame}
\frametitle{Scramble for Africa} 
\begin{figure}
  \includegraphics[keepaspectratio,width=0.9\textwidth]{mp2016_f1.jpeg}
\end{figure}
\end{frame}

\begin{frame}
\frametitle{Data and empirical strategy} 
Data Sources
\begin{itemize}\setlength\itemsep{3pt}
\item Ethnic distribution: Ethnolinguistic map (George Peter Murdock, 1959)
\item Geography: 2000 Digital Chart of the World
  \item Conflict: ACLED, UCDP dataset
  \item Individual-level outcomes: Demographic and Health Surveys
\end{itemize} \medskip
Empirical strategies compare partitioned ethnic groups to non-partitioned ones
\begin{gather*}
\text{Conflict:}\ y_{ec} = \alpha_c + \gamma\text{SPLIT}_{ec}+\phi\text{SPIL}_{ec}+\delta X_{ec}+e_{ec} \\
\text{Indiv:}\ y_{ierc}=\alpha_c + \beta\text{SPLIT}_{e}+ \delta X_{ierc}+\psi Z_{irc}+u_{ierc}
\end{gather*}
\vspace{-2em} 
\begin{itemize}\setlength\itemsep{3pt}
  \item Control for fraction of partitioned adjacent groups in country $c$ $\text{SPIL}_{ec}$
  \item Land and individual characteristics controlled for
  \item For comparability, partitioned and non-partitioned groups need similar characteristics \\ (Only land mass differs)
\end{itemize} \medskip
\end{frame}


\begin{frame}
\frametitle{More conflicts in partitioned groups} 
\begin{figure}
  \includegraphics[keepaspectratio,width=0.9\textwidth]{mp2016_t2.jpeg}
\end{figure}
\end{frame}

\begin{frame}
\frametitle{Partitioned groups experience more discrimination} 
\begin{figure}
  \includegraphics[keepaspectratio,width=0.9\textwidth]{mp2016_t3.jpeg}
\end{figure}
\end{frame}

\begin{frame}
\frametitle{Individuals from partitioned groups are worse-off} 
\begin{figure}
  \includegraphics[keepaspectratio,width=0.67\textwidth]{mp2016_t4.jpeg}
\end{figure}
\end{frame}


\begin{frame}
\frametitle{Discussion and Takeaways} 
Arbitrarily drawn borders $\to$ conflict
\begin{itemize}\setlength\itemsep{3pt}
\item Neighboring countries support co-ethnics on other side of border
\item State-driven conflicts against rebels (of different ethnicities)
\item Political discrimination: Increase participation in ethnic-based civil wars
\end{itemize} \medskip
It also hurts individual well-being
\begin{itemize}\setlength\itemsep{3pt}
\item Fewer assets owned, poor access to public goods, less education
\item Not explained by split country contexts, but due to split ethnicity
\end{itemize} \medskip
\end{frame}

\subsection{Bonnier, Poulsen, Rogall, Stryjan (2020): Preparing for Genocide: Quasi-experimental Evidence from Rwanda [student presentation]}
\begin{transitionframe}
\begin{center}
  \begin{huge}
    \textcolor{NTHU1}{%
Bonnier, Poulsen, Rogall, Stryjan (2020): \\ Preparing for Genocide: Quasi-experimental Evidence from Rwanda [student presentation]
    }
  \end{huge}
\end{center}
\end{transitionframe}


\section{Political motivations behind conflicts}
\subsection{Yanagizawa-Drott (2014): Propaganda and Conflict: Evidence from Rwandan Genocide}
\begin{transitionframe}
\begin{center}
  \begin{huge}
    \textcolor{NTHU1}{%
Yanagizawa-Drott (2014): \\ Propaganda and Conflict: Evidence from Rwandan Genocide
    }
  \end{huge}
\end{center}
\end{transitionframe}

\begin{frame}
\frametitle{Q: What roles do mass media play in conflict and violence?} 
What drives participation into political mass violence?
\begin{itemize}\setlength\itemsep{3pt}
\item Usually sponsored by elites with agendas to eliminate groups considered threatening
\item Reportedly used mass media to justify and induce citizen participation
\begin{itemize}\setlength\itemsep{3pt}
\item Nazis in WW2: Joseph Goebbels as propaganda minister
\item Russia's state-controlled media in its invasion of Ukraine
\end{itemize}
\item So how does these propaganda work?
\end{itemize}\medskip
This paper uses radio campaigns during Rwandan Genocide to answer this question
\begin{itemize}
\item Violence by military, local militias, and ordinary civilians against non-ruling ethnicity
\item Investigates the mechanism of direct persuasion and local spillovers
\begin{itemize}\setlength\itemsep{3pt}
\item Direct persuasion: Convince listeners directly
\item Local spillovers: Spread info and belief via social networks
\end{itemize}
\end{itemize}\medskip
\end{frame}

\begin{frame}
\frametitle{History of Rwandan Genocide} 
Rwandan Genocide in early 1990s
\begin{itemize}\setlength\itemsep{3pt}
\item Hutus majority and Tutsi minority clashed since 1973
\item After assassination of Hutu president (Apr. 1994), extremists began ethnic cleansing
\item State-sponsored militias and civilian-led violence ensued
\item Wave of genocides ended in July 1994: 500,000 Tutsis killed
\end{itemize}\medskip
How was mass media utilized
\begin{itemize} \setlength\itemsep{3pt}
\item Inflammatory messages by RTLM: Run by an extremist Hutu faction (Hutu Power)
\item Called for extermination of Tutsis
\begin{itemize} \setlength\itemsep{3pt}
\item Justified genocide for ``self-defense'' against Tutsis
\item Referred to Tutsis as cockroaches
\item Claimed that Hutus engaged in killings would not be held accountable
\end{itemize}
\end{itemize}\medskip
\end{frame}

\begin{frame}
\frametitle{Conceptual frameworks and mechanisms} 
Direct Effects
\begin{itemize}\setlength\itemsep{3pt}
\item Exposure to media alters behavior by affecting beliefs and preferences
\item Listeners consider RTLM messages as credible: even moderate Hutus were killed
\item But participation in violence is costly: so was RTLM messages persuasive enough?
\end{itemize}\medskip
Spillover Effects
\begin{itemize}\setlength\itemsep{3pt}
\item Social interactions can augment influence of media Exposure
\item Word-of-mouth to nonlisteners, peer pressure, lower cost via collective action 
\item Countervailing effect: Free-rider problem can co-exist in theory
\end{itemize}\medskip
How to distinguish these mechanisms
\begin{itemize}\setlength\itemsep{3pt}
\item Direct effect: Only consider own location
\item Spillover: Also consider characteristics of neighboring villages

\end{itemize}
\end{frame}

\begin{frame}
\frametitle{Data and empirical strategy} 
Sources of key variables
\begin{itemize}\setlength\itemsep{3pt}
\item Violence data from court records (Gacaca courts)
\item Village-level radio data from Bureau of Information and Broadcasting
\item Topography data from Shuttle Radar Topography Mission dataset
\end{itemize}\medskip
Empirical strategy leverages strength of transmission predicted by topography
\begin{itemize}\setlength\itemsep{3pt}
\item i.e. Presence of hills, which is not correlated with other factors affecting conflict
\item Direct and spillovers estimated as follows
\begin{gather*}
\text{Direct:}\ y_{vci}=\beta_v r_{ci}+\pi X_{ci} + \gamma_c + e_{civ} \\
\text{Spillover:}\ y_{vci}=\lambda_{vd} \bar{r}_{dci}+\phi_d \bar{X}_{dci} + \gamma_c + e_{ci}
\end{gather*}
where $\bar{r}$ captures average coverage within $d$km from village $i$ in commune $c$
\end{itemize}\medskip
\end{frame}

\begin{frame}
\frametitle{Distribution of violence and radio} 
\begin{figure}
\begin{subfigure}{0.49\textwidth}
\includegraphics[keepaspectratio,width=\textwidth]{y2014_f1.jpeg}
\caption{Violence distribution}
\end{subfigure}
\begin{subfigure}{0.49\textwidth}
\includegraphics[keepaspectratio,width=\textwidth]{y2014_f2.jpeg}
\caption{Radio distribution}
\end{subfigure}
\end{figure}
\end{frame}

\begin{frame}
\frametitle{Direct effects: Organized violence increases significantly} 
\begin{figure}
\includegraphics[keepaspectratio,width=0.45\textwidth]{y2014_t1.jpeg}
\end{figure}
\end{frame}

\begin{frame}
\frametitle{Spillovers larger than direct effects} 
\begin{figure}
\includegraphics[keepaspectratio,width=0.7\textwidth]{y2014_t2.jpeg}
\end{figure}
\end{frame}

\begin{frame}
\frametitle{Takeaways and discussion} 
Direct effect $<<<<$ Spillover
\begin{itemize}\setlength\itemsep{3pt}
\item Militia presence in \textit{neighboring} villages increased significantly
\item Social interactions do augment influence of direct exposure to media
\item Thus, violence begets violence
\end{itemize}\medskip
Implications
\begin{itemize}\setlength\itemsep{3pt}
\item Need further studies on what type of social interactions matter
 \begin{itemize}\setlength\itemsep{3pt}
\item Collective action change cost-benefit analysis? Social pressure to conform? 
\end{itemize}
\item Dangers of mass media that promotes inflammatory messages
 \begin{itemize}\setlength\itemsep{3pt}
\item Not just violence, but any type of misinformation could have wide-ranging effects
\end{itemize}
\end{itemize}
\end{frame}

\subsection{Condra, Long, Shaver, Wright (2018): The Logic of Insurgent Electoral Violence}
\begin{transitionframe}
\begin{center}
  \begin{huge}
    \textcolor{NTHU1}{%
Condra, Long, Shaver, Wright (2018) \\ The Logic of Insurgent Electoral Violence
    }
  \end{huge}
\end{center}
\end{transitionframe}

\begin{frame}
\frametitle{Q: How are politically-motivated attacks designed?} 
Does elections promote peace or trigger violence?
\begin{itemize}\setlength\itemsep{3pt}
\item Reduce incentives for rebellion and symbolizes state-building and consolidation
  \item Attacks often committed around elections by entities trying to undermine the state
\item Implications for efforts to promote free and fair elections in post-conflict societies
\end{itemize}\medskip
This paper demonstrates logics of insurgents in electoral violence
\begin{itemize}\setlength\itemsep{3pt}
\item Opportunity to undermine legitimacy of the formal government
\item Also cautious enough to minimize civilian targets (sustain civilian support to attackers)
\item Leverages variation in timing and location of violence around election periods
\begin{itemize}\setlength\itemsep{3pt}
  \item Attack earlier in the day than usual?
  \item Damage select routes that voters use to go to polls
\end{itemize}
\end{itemize}
\end{frame}



\begin{frame}
\frametitle{Background and framework} 
Afghanistan in 2001-2017 (before Taliban re-takeover)
\begin{itemize}\setlength\itemsep{3pt}
\item Elections in `04(Pr), `05(Pa), `09(Pr), `10(Pa), `14(Pr) 
\item Taliban actively tried to undermine legitimacy of the Afghan government
\begin{itemize}\setlength\itemsep{3pt}
  \item Targeted candidates, poll booths, security forces, and voters around elections
  \item Qualitative evidence that they calibrated attacks to avoid excessive civilian harm
\end{itemize}
\end{itemize}\medskip
Framework: Competitive governance between insurgents and the state
\begin{itemize}\setlength\itemsep{3pt}
\item Insurgents benefit from violence directed against government
\item But it is undermined if civilians are attacked and consequently
oppose insurgents
\item Electoral violence: Target state legitimacy vs. Harmful when civilians are attacked
\end{itemize} \medskip
Claims to be tested in this paper
\begin{itemize}\setlength\itemsep{3pt}
\item Insurgents time attacks during elections to damage state legitimacy
\item Select location and timing strategically: Use IEDs on key routes and early hours
\end{itemize}
\end{frame}

\begin{frame}
\frametitle{Data and empirical strategy} 
Data Sources
\begin{itemize}\setlength\itemsep{3pt}
\item Attacks (`03-`15): International Security Assistance Force (ISAF)
\item Combined with survey data, population data, voting records, and road network maps
\end{itemize}\medskip
Empirical strategy leverages timing and location of attacks in election vs. non-election days
\begin{itemize}\setlength\itemsep{3pt}
\item Timing: $A_{d_i} = \xi + \sum_{i=1(\neq13)}^{24}\theta_i \text{hour}_{d_i}+ \sum_{i=1(\neq13)}^{24}\alpha_i (\text{hour}_{d_i}\times \text{election}_d)+\psi_d+e_{d_i}$
\item Location: $Y_r = \alpha + \beta_1(\text{Route}_r)+ \beta_2(\text{Violence}_{rw})+\beta_3 X_r + e_r$
\end{itemize} \medskip
Then, the paper estimates the effect of these violence on voting behaviors (IV)
\begin{itemize}\setlength\itemsep{3pt}
\item Morning attacks: $Y_{de}=\alpha +\beta_1 \widehat{\text{Attacks}}_{det}+\beta_2 X_{de}+e_{de}$
\begin{itemize}
\item IV: average of wind conditions in 10:30PM (day before)-4:30AM and 4:30AM-10:30AM
\end{itemize}
\item Location effects: $Y_r=\alpha+\beta_1\widehat{\text{IED}}_r+\beta_2\text{violence}_{rw}+\beta_3 X_r + e_r$
\begin{itemize}
\item IV: Nighttime cloud cover (high cover $\to$ hard to detect $\to$ more attack)
\end{itemize}
\end{itemize} \medskip
\end{frame}

\begin{frame}
\frametitle{Rise of early-morning attacks, but not targeting civilians} 
\begin{figure}
\begin{subfigure}{0.49\textwidth}
\includegraphics[keepaspectratio,width=\textwidth]{clsw2018_f1.jpeg}
\caption{Timing of all attacks}
\end{subfigure}
\begin{subfigure}{0.49\textwidth}
\includegraphics[keepaspectratio,width=\textwidth]{clsw2018_f2.jpeg}
\caption{Timing of attacks causing civilian casualties}
\end{subfigure}
\end{figure}
\end{frame}


\begin{frame}
\frametitle{Rise of IEDs at key roads and population centers, but not more intense} 
\begin{figure}
\includegraphics[keepaspectratio,width=0.8\textwidth]{clsw2018_t1.jpeg}
\end{figure}
\end{frame}

\begin{frame}
\frametitle{Attacks decrease voter turnout} 
\begin{figure}
\begin{subfigure}{0.49\textwidth}
\includegraphics[keepaspectratio,width=\textwidth]{clsw2018_t2.jpeg}
\caption{Effect of morning attacks}
\end{subfigure}
\begin{subfigure}{0.49\textwidth}
\includegraphics[keepaspectratio,width=\textwidth]{clsw2018_t3.jpeg}
\caption{Effect of IEDs}
\end{subfigure}
\end{figure}
\end{frame}

\begin{frame}
\frametitle{Discussion and takeaways} 
Counterfactual exercise to see how alternate violence would have changed outcomes
\begin{itemize}\setlength\itemsep{3pt}
  \item 340-670 IEDs needed to tip election for favor of pro-Taliban parties
  \item Complete elimination of morning attacks $\to$ 12.3 \% $\Uparrow$ in voter turnout
  \item Additional patrol to 1 rotation per month from none $\to$ 5\% $\Downarrow$ Reported frustration
\end{itemize}\medskip
Takeaways
\begin{itemize}\setlength\itemsep{3pt}
  \item Insurgents do compete for governance with the formal state
  \begin{itemize}
    \item This includes other criminal actors like organized criminal groups
  \end{itemize}
  \item Attackers do consider various trade-offs in committing political violence
 \begin{itemize}
    \item Civilian support vs. Undermining state legitimacy
  \end{itemize}
  \item Guidance on designing counterinsurgency efforts: When \& where to invest in security
\end{itemize}
\end{frame}

\subsection{Dube, Naidu (2015): Bases, Bullets, and Ballots: The Effect of US Military Aid on Political Conflict in Colombia [student presentation]}
\begin{transitionframe}
\begin{center}
  \begin{huge}
    \textcolor{NTHU1}{%
Dube, Naidu (2015): \\ Bases, Bullets, and Ballots: The Effect of US Military Aid on Political Conflict in Colombia [student presentation]
    }
  \end{huge}
\end{center}
\end{transitionframe}

\begin{frame}
\frametitle{Takeaways and remaining questions} 
Politics, ethnicity, and institutions affect conflicts
\begin{itemize}\setlength\itemsep{3pt}
\item Ethnicity and other institutional backdrop create differing interests that lead to conflict
\item Politics can propagate incentives to fight and provide opportunities for insurgents
\item Overall, identity has been a driving force behind various conflicts
\end{itemize}\medskip
Unanswered questions
\begin{itemize}\setlength\itemsep{3pt}
\item What to prioritize: Rein in economic incentives or set up inclusive institutions?
\begin{itemize}\setlength\itemsep{3pt}
\item Very complementary: One affects other and vice versa
\item Ongoing debate in the economics of institutions
\end{itemize}
\item Individual participation in these activities?
\begin{itemize}\setlength\itemsep{3pt}
\item We have mostly looked at effects of systematic economic incentives and institutions
\item For individuals, conflicts and crime are very costly. Yet why do the participate?
\end{itemize}
\end{itemize}
\end{frame}

%%%%%%%%%%%
\end{document}


